\documentclass[../../main.tex]{subfiles}
\begin{document}

\begin{flushright}
  \footnotesize\textit{【自動文字起こし・要確認】}
\end{flushright}

% Source: handwritten.pdf 第4問．
{\bf [解]}
$O$を始点とする点$X$の位置ベクトルを$\vec{x}$と表す．題意から
\begin{align}|\vec{a}|=|\vec{b}|=|\vec{c}|.\label{eq:ky2007-4-r}
\end{align}
そして
\begin{align}
  \vec{p}=\dfrac{1}{5}(2\vec{b}+3\vec{a}),\quad
  \vec{q}=\dfrac{1}{5}(2\vec{c}+3\vec{b}),\quad
  \vec{r}=\dfrac{1}{5}(2\vec{a}+3\vec{c}).\label{eq:ky2007-4-div}
\end{align}
$\triangle PQR$の外心が$O$のとき，$|\vec{p}|=|\vec{q}|=|\vec{r}|$となる．
$\triangle ABC$の外接円の半径を1とする．2乗して
$|\vec{p}|^2=|\vec{q}|^2=|\vec{r}|^2=k$とおくと，\cref{eq:ky2007-4-r}から，$k>0$で
\begin{align}
  \begin{cases}
    25k=4+9+12\vec{a}\cdot\vec{b},\\
    25k=13+12\vec{b}\cdot\vec{c},\\
    25k=13+12\vec{c}\cdot\vec{a}.
  \end{cases}
\end{align}
したがって
\begin{align}
  \vec{a}\cdot\vec{b}=\vec{b}\cdot\vec{c}=\vec{c}\cdot\vec{a}=\dfrac{25k-13}{12}.
  \label{eq:ky2007-4-dot}
\end{align}
ここで対称性から$O$を原点とする座標平面で$A(1,0)$，
$B(\cos\alpha,\sin\alpha)$，$C(\cos\beta,\sin\beta)$とおく．
$0<\alpha<\beta<2\pi$であり，\cref{eq:ky2007-4-dot}から
\begin{align}\cos\alpha=\cos\beta=\cos(\beta-\alpha).
\end{align}
$0<\alpha<\beta<2\pi$から，左側の等号から$\beta=2\pi-\alpha$とかける．
\begin{align}0<\alpha<2\pi-\alpha<2\pi,\qquad0<\alpha<\pi.\label{eq:ky2007-4-angle}
\end{align}
したがって右側の等号から
\begin{align}
  \cos\alpha=\cos(2\pi-2\alpha)=\cos2\alpha=2\cos^2\alpha-1,\\
  (2\cos\alpha+1)(\cos\alpha-1)=0.
\end{align}
\cref{eq:ky2007-4-angle}から$\alpha=\dfrac{2\pi}{3}$，$\beta=2\pi-\alpha=\dfrac{4\pi}{3}$．
したがって$\triangle ABC$は正三角形．

\medskip\noindent{\bf [別解]}
\cref{eq:ky2007-4-dot}以外に，\cref{eq:ky2007-4-div}から
\begin{align}|AB|^2=|\vec{b}-\vec{a}|^2
\end{align}
などとして，$\vec{a}\cdot\vec{b}=\vec{b}\cdot\vec{c}=\vec{c}\cdot\vec{a}$から
$|AB|=|BC|=|CA|$となる．この方が早い．

\end{document}
